% Generated draft: requires mathematical and source review.
\documentclass[11pt]{article}
\usepackage[margin=25mm]{geometry}
\usepackage{fontspec}
\setmainfont{Latin Modern Roman}
\usepackage{amsmath,amssymb,mathtools,unicode-math}
\setmathfont{Latin Modern Math}
\usepackage{enumitem,needspace,xurl,hyperref}
\setlength{\parindent}{0pt}
\setlength{\parskip}{5pt}
\setlength{\emergencystretch}{3em}
\newcommand{\N}{\mathbb{N}}
\newcommand{\Z}{\mathbb{Z}}
\newcommand{\Q}{\mathbb{Q}}
\newcommand{\R}{\mathbb{R}}
\newcommand{\C}{\mathbb{C}}
\newcommand{\F}{\mathbb{F}}
\newcommand{\A}{\mathbb{A}}
\newcommand{\PP}{\mathbb{P}}
\newcommand{\E}{\mathbb{E}}
\newcommand{\eps}{\varepsilon}
\newcommand{\dd}{\,\mathrm{d}}
\begin{document}
\section*{UnsolvedMath ID: 30003594}
\textbf{Source problem code:} catalogue entry
\subsection{Definitions and mathematical statement}
Let \(X\) be a smooth complex projective variety and let \(V\subset X\) be a subvariety. Let \(D\) be a divisor with coefficients in \(\Q\) or \(\R\). Since \(X\) is smooth, such divisors can be viewed as \(\Q\)-Cartier or \(\R\)-Cartier divisors, respectively. Two real divisors \(D_1,D_2\) are numerically equivalent if
\[
D_1\cdot C=D_2\cdot C
\qquad\text{for every integral curve }C\subset X.
\]
Write \(N^1(X)_{\R}\) for the real vector space of real divisor classes modulo numerical equivalence and \([D]\) for the numerical class of \(D\). The pseudoeffective cone \(\overline{\operatorname{Eff}}^{\,1}(X)\) is the closure in \(N^1(X)_{\R}\) of the cone generated by effective divisors. The assumption that \(D\) is big means
\[
[D]\in\operatorname{int}_{N^1(X)_{\R}}\!\left(\overline{\operatorname{Eff}}^{\,1}(X)\right).
\]
\textbf{Source problem, with unresolved scope.} For every such \(X\), \(V\), and big \(D\), determine the geometry of \([D]\) relative to \(V\); the supplied edited formulation describes this as determining its positivity and birational geometry relative to \(V\).

The source does not define what constitutes this determination, which relative positivity properties or birational constructions are intended, or what data a characterization may use. Thus this is a research question rather than a specified proposition with a definite conclusion. No additional hypothesis on \(V\), including smoothness, dimension, or position relative to the base loci of \(D\), is supplied. Nor is a convention about irreducibility of \(V\) explicitly stated. These qualifications remain unresolved; no particular interpretation of geometry relative to \(V\) is selected here.
\subsection{Short English statement}
Given a smooth complex projective variety \(X\), a subvariety \(V\subset X\), and a big rational or real divisor \(D\), the source asks how to describe the geometry of the numerical class \([D]\) relative to \(V\). The edited wording emphasizes positivity and birational geometry, but the intended properties and required form of the answer are unspecified. This statement-only draft requires human mathematical review.
\subsection{Sources}
\begin{itemize}
\item Local UnsolvedMath registry, permanent ID \(30003594\), title: \emph{Big Divisor Geometry Relative to Subvarieties}; collection: Oberwolfach Reports Open Problems; exact source code: \texttt{catalogue entry}. The supplied record gives proposed year 2017 and names no proposer.
\item Supplied report citation: \emph{Mini-Workshop: Positivity in Higher-dimensional Geometry: Higher-codimensional Cycles and Newton--Okounkov Bodies} (2018), Oberwolfach Reports, pp.\ 2631--2657, DOI \texttt{10.4171/owr/2017/43}. Supplied source URL: \url{https://doi.org/10.4171/owr/2017/43}. The record describes Oberwolfach Reports as documenting mathematical workshops and their open-problem discussions. The report citation is supplied provenance and has not been independently inspected here.
\item Supplied literature source, labeled ``Lozovanu--Kuronya, continuity and convexity'': \url{https://ems.press/content/serial-article-files/46714}.
\item Supplied literature source, labeled ``Kuronya, asymptotic invariants of base loci'': \url{https://arxiv.org/abs/math/0308116}.
\item Supplied literature source, labeled ``original OWR report'': \url{https://doi.org/10.4171/OWR/2017/43}.
\item Literature assessment dated \textbf{2026-08-21}: the supplied curation labels the problem open and states that intersection numbers, restricted volumes, and positivity notions provide partial information, while the requested general geometric determination has no settled answer. This is supplied dated curation derived from the OWR material, not an independently verified literature assessment. The record itself requests independent verification of time-sensitive status and cited claims.
\item The supplied statement status is \texttt{corrected\_verified}, accompanied by the explanation that notation, typography, and sentence structure were repaired while preserving mathematical content. This is a supplied editorial label, not independent mathematical verification of this draft or its interpretation.
\end{itemize}
\end{document}
